\subsection{A proposed alternating optimization of the two forests}
\label{sec:paired-forests}

The implemented weighted dual-tree choice optimizes the retained literal
attaching-curve work after fixing the primal forest. There is a symmetric,
currently unimplemented optimization of the primal forest. The following
statement concerns two finite graphs and an incidence matrix; its proof
requires no additional claim about transporting boundary markings.

Let $G_{01}$ be a finite multigraph with edge set $\mathcal H$, and let
$G_{23}$ be a connected finite multigraph with distinguished root and edge
set $\mathcal R$. Loops are allowed in both graphs but cannot lie in a
forest. Let $N=(n_{rh})_{r\in\mathcal R,h\in\mathcal H}$ have nonnegative
integer entries. In the intended source application, $n_{rh}$ is the
number of literal passages of source two-handle $r$ through source
one-handle $h$. Let $T$ span every connected component of $G_{01}$ and let
$D$ be a spanning tree of $G_{23}$. Define
\[
 C(T,D)=\sum_{r\notin D}\sum_{h\notin T}n_{rh},
 \qquad M=\sum_{r,h}n_{rh}.
\]
Thus $C$ counts retained source passages after deleting the selected
source two-handles and ignoring the primal-tree one-handles. It uses raw
occurrences, without cancellation of inverse letters or any prediction of
later handle slides.

\begin{proposition}[Alternating exact forest responses]
\label{prop:paired-forests}
For fixed $T$, a maximum-weight spanning tree of $G_{23}$ with edge
weights
\[
 w_T(r)=\sum_{h\notin T}n_{rh}
\]
minimizes $C(T,D)$. For fixed $D$, a maximum-weight spanning forest of
$G_{01}$, spanning each of its components, with edge weights
\[
 v_D(h)=\sum_{r\notin D}n_{rh}
\]
minimizes $C(T,D)$ over $T$.

Starting from any feasible pair, alternate these two exact responses,
accepting a changed forest only when it strictly decreases $C$. Stop after
a complete sweep in which neither response gives a strict decrease. There
are at most $C(T_0,D_0)\le M$ accepted updates. At termination, changing
either forest alone cannot lower $C$.

Suppose the graphs have total combinatorial size $O(t)$ and the matrix is
provided as literal occurrence lists of total length $M$. Then this
procedure uses
\[
 O\bigl((M+1)(M+t\log(t+2))\bigr)
\]
word operations and $O(M+t)$ working storage. In particular, for source
incidence of size $M=O(t)$, the bound is $O(t^2\log(t+2))$ after source
construction and authentication.
\end{proposition}

\begin{proof}
For fixed $T$,
\[
 C(T,D)=\sum_{r\in\mathcal R}w_T(r)-\sum_{r\in D}w_T(r).
\]
The first term is constant, proving the dual assertion. For fixed $D$,
\[
 C(T,D)=\sum_{h\in\mathcal H}v_D(h)-\sum_{h\in T}v_D(h),
\]
which proves the primal assertion component by component. These formulas
include all loop weights in the constant terms; a loop is never silently
treated as a cancellable forest edge.

The objective is a nonnegative integer, so every accepted update lowers it
by at least one. Every nonterminal full sweep contains an accepted update.
There are consequently at most $M+1$ sweeps, including the final sweep.
If an optimal response has the current objective value, retaining the
current forest is essential to the specified stopping rule: zero-gain
changes are not performed. Thus, in the final unchanged sweep, both
current forests individually attain their conditional minimum.

One pass through the occurrence lists computes the required weights.
Sorting weighted graph edges and running the usual maximum-spanning-forest
algorithm costs $O(t\log(t+2))$ word operations per response. There are at
most two responses per sweep. All weights lie between zero and $M$, so
$O(\log(M+t+2))$-bit words suffice for weights and identifiers. This proves
the complexity and storage bounds.
\end{proof}

Conditional optimality can be strictly weaker than joint optimality even
for two-edge graphs. Give each graph two vertices joined by two parallel
edges and take
\[
 N=\begin{pmatrix}1&2\\2&0\end{pmatrix}.
\]
Deleting row two and column two leaves cost one. Changing either deletion
alone raises the cost to two, so both responses retain this pair. Deleting
row one and column one together instead leaves cost zero. This is an
abstract incidence example, without a claimed geometric realization.

This is a proposal for a producer heuristic, not a modification of the
runtime delivered here. A final chosen dual tree would still require the
independent source-side replay, and a chosen primal forest would still
need to be recorded in the applicable downstream certificate. The graph
objective does not establish transport of an arbitrary boundary pattern
through primal contractions. Nor does conditional optimality prove global
optimality over pairs, existence of a complete raw cancellation schedule,
or completeness of the recognition algorithm. The strict-update bound is
polynomial for the stated literal source representation; for an arbitrary
matrix of huge binary-encoded multiplicities, a bound proportional to the
numerical value $M$ would instead be pseudo-polynomial.
